% Options for packages loaded elsewhere
\PassOptionsToPackage{unicode}{hyperref}
\PassOptionsToPackage{hyphens}{url}
\documentclass[
  11pt,
]{article}
\usepackage{xcolor}
\usepackage[margin=1in]{geometry}
\usepackage{amsmath,amssymb}
\setcounter{secnumdepth}{-\maxdimen} % remove section numbering
\usepackage{iftex}
\ifPDFTeX
  \usepackage[T1]{fontenc}
  \usepackage[utf8]{inputenc}
  \usepackage{textcomp} % provide euro and other symbols
\else % if luatex or xetex
  \usepackage{unicode-math} % this also loads fontspec
  \defaultfontfeatures{Scale=MatchLowercase}
  \defaultfontfeatures[\rmfamily]{Ligatures=TeX,Scale=1}
\fi
\usepackage{lmodern}
\ifPDFTeX\else
  % xetex/luatex font selection
\fi
% Use upquote if available, for straight quotes in verbatim environments
\IfFileExists{upquote.sty}{\usepackage{upquote}}{}
\IfFileExists{microtype.sty}{% use microtype if available
  \usepackage[]{microtype}
  \UseMicrotypeSet[protrusion]{basicmath} % disable protrusion for tt fonts
}{}
\makeatletter
\@ifundefined{KOMAClassName}{% if non-KOMA class
  \IfFileExists{parskip.sty}{%
    \usepackage{parskip}
  }{% else
    \setlength{\parindent}{0pt}
    \setlength{\parskip}{6pt plus 2pt minus 1pt}}
}{% if KOMA class
  \KOMAoptions{parskip=half}}
\makeatother
\usepackage{color}
\usepackage{fancyvrb}
\newcommand{\VerbBar}{|}
\newcommand{\VERB}{\Verb[commandchars=\\\{\}]}
\DefineVerbatimEnvironment{Highlighting}{Verbatim}{commandchars=\\\{\}}
% Add ',fontsize=\small' for more characters per line
\newenvironment{Shaded}{}{}
\newcommand{\AlertTok}[1]{\textcolor[rgb]{1.00,0.00,0.00}{\textbf{#1}}}
\newcommand{\AnnotationTok}[1]{\textcolor[rgb]{0.38,0.63,0.69}{\textbf{\textit{#1}}}}
\newcommand{\AttributeTok}[1]{\textcolor[rgb]{0.49,0.56,0.16}{#1}}
\newcommand{\BaseNTok}[1]{\textcolor[rgb]{0.25,0.63,0.44}{#1}}
\newcommand{\BuiltInTok}[1]{\textcolor[rgb]{0.00,0.50,0.00}{#1}}
\newcommand{\CharTok}[1]{\textcolor[rgb]{0.25,0.44,0.63}{#1}}
\newcommand{\CommentTok}[1]{\textcolor[rgb]{0.38,0.63,0.69}{\textit{#1}}}
\newcommand{\CommentVarTok}[1]{\textcolor[rgb]{0.38,0.63,0.69}{\textbf{\textit{#1}}}}
\newcommand{\ConstantTok}[1]{\textcolor[rgb]{0.53,0.00,0.00}{#1}}
\newcommand{\ControlFlowTok}[1]{\textcolor[rgb]{0.00,0.44,0.13}{\textbf{#1}}}
\newcommand{\DataTypeTok}[1]{\textcolor[rgb]{0.56,0.13,0.00}{#1}}
\newcommand{\DecValTok}[1]{\textcolor[rgb]{0.25,0.63,0.44}{#1}}
\newcommand{\DocumentationTok}[1]{\textcolor[rgb]{0.73,0.13,0.13}{\textit{#1}}}
\newcommand{\ErrorTok}[1]{\textcolor[rgb]{1.00,0.00,0.00}{\textbf{#1}}}
\newcommand{\ExtensionTok}[1]{#1}
\newcommand{\FloatTok}[1]{\textcolor[rgb]{0.25,0.63,0.44}{#1}}
\newcommand{\FunctionTok}[1]{\textcolor[rgb]{0.02,0.16,0.49}{#1}}
\newcommand{\ImportTok}[1]{\textcolor[rgb]{0.00,0.50,0.00}{\textbf{#1}}}
\newcommand{\InformationTok}[1]{\textcolor[rgb]{0.38,0.63,0.69}{\textbf{\textit{#1}}}}
\newcommand{\KeywordTok}[1]{\textcolor[rgb]{0.00,0.44,0.13}{\textbf{#1}}}
\newcommand{\NormalTok}[1]{#1}
\newcommand{\OperatorTok}[1]{\textcolor[rgb]{0.40,0.40,0.40}{#1}}
\newcommand{\OtherTok}[1]{\textcolor[rgb]{0.00,0.44,0.13}{#1}}
\newcommand{\PreprocessorTok}[1]{\textcolor[rgb]{0.74,0.48,0.00}{#1}}
\newcommand{\RegionMarkerTok}[1]{#1}
\newcommand{\SpecialCharTok}[1]{\textcolor[rgb]{0.25,0.44,0.63}{#1}}
\newcommand{\SpecialStringTok}[1]{\textcolor[rgb]{0.73,0.40,0.53}{#1}}
\newcommand{\StringTok}[1]{\textcolor[rgb]{0.25,0.44,0.63}{#1}}
\newcommand{\VariableTok}[1]{\textcolor[rgb]{0.10,0.09,0.49}{#1}}
\newcommand{\VerbatimStringTok}[1]{\textcolor[rgb]{0.25,0.44,0.63}{#1}}
\newcommand{\WarningTok}[1]{\textcolor[rgb]{0.38,0.63,0.69}{\textbf{\textit{#1}}}}
\setlength{\emergencystretch}{3em} % prevent overfull lines
\providecommand{\tightlist}{%
  \setlength{\itemsep}{0pt}\setlength{\parskip}{0pt}}
\usepackage{bookmark}
\IfFileExists{xurl.sty}{\usepackage{xurl}}{} % add URL line breaks if available
\urlstyle{same}
\hypersetup{
  hidelinks,
  pdfcreator={LaTeX via pandoc}}

\author{}
\date{}

\begin{document}

\section{Final Exam Question Bank}\label{final-exam-question-bank}

This file contains the question bank for the final exam.

There are 5 types of questions, which are worth the following points:

\begin{itemize}
\tightlist
\item
  select all that are true = 2 points
\item
  short answer = 2 points
\item
  predict the output = 3 points
\item
  short coding challenge = 4 points
\item
  long coding challenge = 8 points
\end{itemize}

The final exam will be organized into five parts:

\begin{itemize}
\tightlist
\item
  Part 1: 8 points, 4 questions, 2 points each -- select all that are
  true
\item
  Part 2: 8 points, 4 questions, 2 points each -- short answer questions
\item
  Part 3: 12 points, 4 questions, 3 points each -- predict the output
  and explain why
\item
  Part 4: 8 points, 2 questions, 4 points each -- short coding
  challenges
\item
  Part 5: 16 points, 2 questions, 8 points each -- longer coding
  challenges
\end{itemize}

The final exam will be worth 52 points and the questions will be sampled
from this question bank.

You'll have 2 hours (120 minutes) to complete the exam.

The exam is closed notes, no electronic devices, but you will be given
the following \url{final_exam_reference_material.md} at the exam.

\subsection{Select All That Are True
(mcq\_multi)}\label{select-all-that-are-true-mcq_multi}

Select all statements that are true about ownership and moves in Rust:

\begin{itemize}
\tightlist
\item[$\square$]
  A. Assigning a \texttt{String} to another variable moves ownership;
  the original variable is no longer valid.
\item[$\square$]
  B. Passing a \texttt{\&String} (a shared reference) to a function
  moves ownership of the \texttt{String} into the function.
\item[$\square$]
  C. Assigning an \texttt{i32} to another variable moves ownership; the
  original variable is no longer valid.
\item[$\square$]
  D. A value is dropped when its owner goes out of scope.
\end{itemize}

\begin{center}\rule{0.5\linewidth}{0.5pt}\end{center}

Select all statements that are true about borrowing rules at a single
point in a program:

\begin{itemize}
\tightlist
\item[$\square$]
  A. You can have any number of immutable references (\texttt{\&T})
  simultaneously.
\item[$\square$]
  B. You can have one mutable reference (\texttt{\&mut\ T}) and one
  immutable reference simultaneously.
\item[$\square$]
  C. You can have at most one mutable reference at a time.
\item[$\square$]
  D. A mutable reference \texttt{\&mut\ T} gives you read-only access to
  the value it points at.
\end{itemize}

\begin{center}\rule{0.5\linewidth}{0.5pt}\end{center}

Select all statements that are true about traits in Rust:

\begin{itemize}
\tightlist
\item[$\square$]
  A. A trait defines a set of method signatures (and optionally default
  implementations) that types can implement.
\item[$\square$]
  B. You can implement a trait you did not define on a type you did not
  define.
\item[$\square$]
  C. A single type can implement multiple traits.
\item[$\square$]
  D. A trait method can have a default implementation that implementors
  may override.
\end{itemize}

\begin{center}\rule{0.5\linewidth}{0.5pt}\end{center}

Select all statements that are true about \texttt{\#{[}derive(...){]}}
in Rust:

\begin{itemize}
\tightlist
\item[$\square$]
  A. \texttt{\#{[}derive(Debug){]}} lets you print a value with
  \texttt{\{:?\}}.
\item[$\square$]
  B. To derive \texttt{Copy}, the type must also derive (or implement)
  \texttt{Clone}.
\item[$\square$]
  C. You can derive \texttt{Copy} on a struct that has a \texttt{String}
  field.
\item[$\square$]
  D. To use a struct as a \texttt{HashMap} key, the struct typically
  needs to derive \texttt{Eq} and \texttt{Hash}.
\end{itemize}

\begin{center}\rule{0.5\linewidth}{0.5pt}\end{center}

Select all statements that are true about lifetimes in Rust:

\begin{itemize}
\tightlist
\item[$\square$]
  A. Lifetimes describe how long references are valid and are checked at
  compile time.
\item[$\square$]
  B. Lifetime annotations change how long a value lives at runtime.
\item[$\square$]
  C. A function signature like
  \texttt{fn\ longest\textless{}\textquotesingle{}a\textgreater{}(x:\ \&\textquotesingle{}a\ str,\ y:\ \&\textquotesingle{}a\ str)\ -\textgreater{}\ \&\textquotesingle{}a\ str}
  says the returned reference is valid for at least as long as both
  inputs.
\item[$\square$]
  D. Every function that takes a reference parameter must include an
  explicit lifetime annotation written out by the programmer.
\end{itemize}

\begin{center}\rule{0.5\linewidth}{0.5pt}\end{center}

Select all statements that are true about generics in Rust:

\begin{itemize}
\tightlist
\item[$\square$]
  A.
  \texttt{fn\ largest\textless{}T:\ PartialOrd\textgreater{}(list:\ \&{[}T{]})\ -\textgreater{}\ \&T}
  works for any type \texttt{T} that can be ordered.
\item[$\square$]
  B. Generics are resolved at runtime through dynamic dispatch by
  default.
\item[$\square$]
  C. A trait bound like \texttt{T:\ Clone\ +\ Debug} means that
  \texttt{T} must implement at least one of \texttt{Clone} or
  \texttt{Debug}.
\item[$\square$]
  D. Generic functions always rely on runtime type information (RTTI) to
  decide which type-specific code to execute.
\end{itemize}

\begin{center}\rule{0.5\linewidth}{0.5pt}\end{center}

Select all statements that are true about iterators in Rust:

\begin{itemize}
\tightlist
\item[$\square$]
  A. Calling \texttt{.map(...)} on an iterator does the work
  immediately.
\item[$\square$]
  B. Iterator adaptors like \texttt{filter}, \texttt{map}, and
  \texttt{take} are lazy --- no work happens until a consuming method
  runs.
\item[$\square$]
  C. \texttt{.collect()} is a consuming method that drives the iterator
  to completion.
\item[$\square$]
  D. \texttt{.sum()} is a consuming method that returns the total of all
  elements.
\end{itemize}

\begin{center}\rule{0.5\linewidth}{0.5pt}\end{center}

Select all statements that are true about closures in Rust:

\begin{itemize}
\tightlist
\item[$\square$]
  A. A closure can capture variables from its enclosing scope by
  reference, by mutable reference, or by move.
\item[$\square$]
  B. The \texttt{move} keyword forces the closure to take ownership of
  captured variables.
\item[$\square$]
  C. A closure and a regular \texttt{fn} function have exactly the same
  type in Rust.
\item[$\square$]
  D. Closures are commonly passed to iterator methods such as
  \texttt{map}, \texttt{filter}, and \texttt{fold}.
\end{itemize}

\begin{center}\rule{0.5\linewidth}{0.5pt}\end{center}

Select all statements that are true about
\texttt{Option\textless{}T\textgreater{}} and
\texttt{Result\textless{}T,\ E\textgreater{}}:

\begin{itemize}
\tightlist
\item[$\square$]
  A. \texttt{Option\textless{}T\textgreater{}} represents a value that
  may or may not be present.
\item[$\square$]
  B. \texttt{Result\textless{}T,\ E\textgreater{}} represents either a
  success value of type \texttt{T} or an error of type \texttt{E}.
\item[$\square$]
  C. Calling \texttt{.unwrap()} on \texttt{None} returns the default
  value of type \texttt{T} rather than panicking.
\item[$\square$]
  D. \texttt{match} cannot be used to destructure \texttt{Option} or
  \texttt{Result} values.
\end{itemize}

\begin{center}\rule{0.5\linewidth}{0.5pt}\end{center}

Select all statements that are true about
\texttt{HashMap\textless{}K,\ V\textgreater{}}:

\begin{itemize}
\tightlist
\item[$\square$]
  A. Keys must implement \texttt{Eq} and \texttt{Hash}.
\item[$\square$]
  B. Iteration order is guaranteed to match insertion order.
\item[$\square$]
  C. Average \texttt{insert}, \texttt{get}, and \texttt{remove} are
  O(1).
\item[$\square$]
  D. Calling \texttt{insert} on an existing key replaces the old value
  and returns \texttt{Some(old\_value)}.
\end{itemize}

\begin{center}\rule{0.5\linewidth}{0.5pt}\end{center}

Select all statements that are true about
\texttt{BTreeMap\textless{}K,\ V\textgreater{}}:

\begin{itemize}
\tightlist
\item[$\square$]
  A. Keys must implement \texttt{Ord}.
\item[$\square$]
  B. Iteration visits entries in insertion order.
\item[$\square$]
  C. \texttt{insert}, \texttt{get}, and \texttt{remove} are O(1)
  average-case, just like \texttt{HashMap}.
\item[$\square$]
  D. \texttt{first\_key\_value()} returns the entry that was inserted
  first in time.
\end{itemize}

\begin{center}\rule{0.5\linewidth}{0.5pt}\end{center}

Select all statements that are true about
\texttt{HashSet\textless{}T\textgreater{}}:

\begin{itemize}
\tightlist
\item[$\square$]
  A. A \texttt{HashSet} stores unique values --- inserting an existing
  value has no effect on set membership.
\item[$\square$]
  B. \texttt{insert} returns \texttt{true} if the value was newly
  inserted and \texttt{false} if it was already present.
\item[$\square$]
  C. \texttt{HashSet} iteration is guaranteed to visit elements in the
  order they were inserted.
\item[$\square$]
  D. \texttt{HashSet} keeps its elements in sorted order.
\end{itemize}

\begin{center}\rule{0.5\linewidth}{0.5pt}\end{center}

Select all statements that are true about
\texttt{BinaryHeap\textless{}T\textgreater{}} in Rust:

\begin{itemize}
\tightlist
\item[$\square$]
  A. By default, \texttt{BinaryHeap\textless{}T\textgreater{}} is a
  max-heap: \texttt{pop()} returns the largest element.
\item[$\square$]
  B. Elements only need to implement \texttt{PartialOrd}, so a
  \texttt{BinaryHeap\textless{}f64\textgreater{}} compiles and works
  as-is.
\item[$\square$]
  C. Iterating with \texttt{.iter()} returns elements in sorted order.
\item[$\square$]
  D. \texttt{BinaryHeap::peek()} removes and returns the largest
  element.
\end{itemize}

\begin{center}\rule{0.5\linewidth}{0.5pt}\end{center}

Select all statements that are true about stacks and queues in Rust:

\begin{itemize}
\tightlist
\item[$\square$]
  A. A \texttt{Vec\textless{}T\textgreater{}} can be used as a stack via
  \texttt{push} / \texttt{pop}, with both operations running in
  amortized O(1).
\item[$\square$]
  B. Stacks are LIFO (last-in, first-out) and queues are FIFO (first-in,
  first-out).
\item[$\square$]
  C. \texttt{VecDeque\textless{}T\textgreater{}} supports efficient
  \texttt{push\_back} and \texttt{pop\_front}, making it a natural FIFO
  queue.
\item[$\square$]
  D. \texttt{Vec::remove(0)} is an O(1) operation.
\end{itemize}

\begin{center}\rule{0.5\linewidth}{0.5pt}\end{center}

Select all statements that are true about the (average-case) complexity
of common operations:

\begin{itemize}
\tightlist
\item[$\square$]
  A. Inserting at the front of a \texttt{Vec\textless{}T\textgreater{}}
  with \texttt{vec.insert(0,\ x)} is O(1).
\item[$\square$]
  B. \texttt{HashMap::get} is O(1) on average.
\item[$\square$]
  C. \texttt{BTreeMap::get} is O(1) on average, just like
  \texttt{HashMap::get}.
\item[$\square$]
  D. Looking up an element by value in an unsorted
  \texttt{Vec\textless{}T\textgreater{}} with \texttt{.contains(\&x)} is
  O(log n).
\end{itemize}

\begin{center}\rule{0.5\linewidth}{0.5pt}\end{center}

Select all statements that are true about the \texttt{Entry} API on
\texttt{HashMap} and \texttt{BTreeMap}:

\begin{itemize}
\tightlist
\item[$\square$]
  A. \texttt{map.entry(key).or\_insert(default)} inserts
  \texttt{default} only if the key is absent, and returns a mutable
  reference to the value.
\item[$\square$]
  B. \texttt{map.entry(key).or\_insert(default)} always overwrites the
  existing value when the key is already present.
\item[$\square$]
  C. The \texttt{Entry} API is available only on \texttt{HashMap};
  \texttt{BTreeMap} does not provide an \texttt{entry} method.
\item[$\square$]
  D. Internally, \texttt{map.entry(key).or\_insert(v)} performs two
  separate lookups (one to check for the key, one to insert), so it is
  slower than a manual \texttt{get} + \texttt{insert}.
\end{itemize}

\begin{center}\rule{0.5\linewidth}{0.5pt}\end{center}

\subsection{Short Answer
(short\_answer)}\label{short-answer-short_answer}

\textbf{Why does Rust have an ownership system? Name one problem it
prevents at compile time that languages like C or Python handle
differently.}

\begin{center}\rule{0.5\linewidth}{0.5pt}\end{center}

\textbf{What is a trait in Rust? Give one concrete example of a
standard-library trait and describe what it lets you do.}

\begin{center}\rule{0.5\linewidth}{0.5pt}\end{center}

\textbf{What does a lifetime annotation like
\texttt{\textquotesingle{}a} in
\texttt{fn\ longest\textless{}\textquotesingle{}a\textgreater{}(x:\ \&\textquotesingle{}a\ str,\ y:\ \&\textquotesingle{}a\ str)\ -\textgreater{}\ \&\textquotesingle{}a\ str}
tell the Rust compiler?}

\begin{center}\rule{0.5\linewidth}{0.5pt}\end{center}

\textbf{Why are generics useful in Rust? Give one example.}

\begin{center}\rule{0.5\linewidth}{0.5pt}\end{center}

\textbf{What does it mean that iterator adaptors in Rust are ``lazy''?
Give an example of an adaptor and an example of a consumer.}

\begin{center}\rule{0.5\linewidth}{0.5pt}\end{center}

\textbf{How does a closure in Rust differ from a regular \texttt{fn}
function?}

\begin{center}\rule{0.5\linewidth}{0.5pt}\end{center}

\textbf{What is the difference between
\texttt{Option\textless{}T\textgreater{}} and
\texttt{Result\textless{}T,\ E\textgreater{}}, and when would you reach
for each?}

\begin{center}\rule{0.5\linewidth}{0.5pt}\end{center}

\textbf{What is the role of the hash function when using a
\texttt{HashMap\textless{}K,\ V\textgreater{}}, and why must keys
implement \texttt{Eq} as well as \texttt{Hash}?}

\begin{center}\rule{0.5\linewidth}{0.5pt}\end{center}

\textbf{When would you choose
\texttt{BTreeMap\textless{}K,\ V\textgreater{}} over
\texttt{HashMap\textless{}K,\ V\textgreater{}}? Name one specific
advantage.}

\begin{center}\rule{0.5\linewidth}{0.5pt}\end{center}

\textbf{Give one concrete use case where a
\texttt{HashSet\textless{}T\textgreater{}} is a better choice than a
\texttt{Vec\textless{}T\textgreater{}}, and briefly explain why.}

\begin{center}\rule{0.5\linewidth}{0.5pt}\end{center}

\textbf{What is \texttt{BinaryHeap\textless{}T\textgreater{}} useful
for? Describe one typical use case and the order in which \texttt{pop()}
returns elements.}

\begin{center}\rule{0.5\linewidth}{0.5pt}\end{center}

\textbf{Explain the difference between a stack and a queue. Which Rust
standard-library type is a good fit for each?}

\begin{center}\rule{0.5\linewidth}{0.5pt}\end{center}

\textbf{Describe the difference between O(1), O(n), and O(log n) time
complexity in plain language. Give one Rust collection operation for
each.}

\begin{center}\rule{0.5\linewidth}{0.5pt}\end{center}

\textbf{What problem does the \texttt{Entry} API on \texttt{HashMap} /
\texttt{BTreeMap} solve, compared with using \texttt{contains\_key}
followed by \texttt{insert}?}

\begin{center}\rule{0.5\linewidth}{0.5pt}\end{center}

\textbf{Name three commonly derived traits in Rust and briefly describe
what \texttt{\#{[}derive(...){]}} gives you for each.}

\begin{center}\rule{0.5\linewidth}{0.5pt}\end{center}

\textbf{Give one reason why \texttt{Vec\textless{}T\textgreater{}} is
usually preferred over \texttt{LinkedList\textless{}T\textgreater{}} in
Rust, even when you need to frequently grow a collection.}

\begin{center}\rule{0.5\linewidth}{0.5pt}\end{center}

\subsection{Predict the Output
(predict\_output)}\label{predict-the-output-predict_output}

\textbf{What is the output of this program?}

\begin{Shaded}
\begin{Highlighting}[]
\KeywordTok{fn}\NormalTok{ main() }\OperatorTok{\{}
    \KeywordTok{let}\NormalTok{ x }\OperatorTok{=} \DecValTok{5}\OperatorTok{;}
    \KeywordTok{let}\NormalTok{ y }\OperatorTok{=}\NormalTok{ x}\OperatorTok{;}
    \KeywordTok{let}\NormalTok{ x }\OperatorTok{=}\NormalTok{ x }\OperatorTok{+} \DecValTok{1}\OperatorTok{;}
    \PreprocessorTok{println!}\NormalTok{(}\StringTok{"\{\} \{\}"}\OperatorTok{,}\NormalTok{ x}\OperatorTok{,}\NormalTok{ y)}\OperatorTok{;}
\OperatorTok{\}}
\end{Highlighting}
\end{Shaded}

\begin{center}\rule{0.5\linewidth}{0.5pt}\end{center}

\textbf{What is the output of this program?}

\begin{Shaded}
\begin{Highlighting}[]
\KeywordTok{fn}\NormalTok{ double}\OperatorTok{\textless{}}\NormalTok{T}\OperatorTok{:} \PreprocessorTok{std::ops::}\BuiltInTok{Add}\OperatorTok{\textless{}}\NormalTok{Output }\OperatorTok{=}\NormalTok{ T}\OperatorTok{\textgreater{}} \OperatorTok{+} \BuiltInTok{Copy}\OperatorTok{\textgreater{}}\NormalTok{(x}\OperatorTok{:}\NormalTok{ T) }\OperatorTok{{-}\textgreater{}}\NormalTok{ T }\OperatorTok{\{}
\NormalTok{    x }\OperatorTok{+}\NormalTok{ x}
\OperatorTok{\}}

\KeywordTok{fn}\NormalTok{ main() }\OperatorTok{\{}
    \PreprocessorTok{println!}\NormalTok{(}\StringTok{"\{\} \{\}"}\OperatorTok{,}\NormalTok{ double(}\DecValTok{3}\NormalTok{)}\OperatorTok{,}\NormalTok{ double(}\DecValTok{2.5}\NormalTok{))}\OperatorTok{;}
\OperatorTok{\}}
\end{Highlighting}
\end{Shaded}

\begin{center}\rule{0.5\linewidth}{0.5pt}\end{center}

\textbf{What is the output of this program?}

\begin{Shaded}
\begin{Highlighting}[]
\KeywordTok{fn}\NormalTok{ main() }\OperatorTok{\{}
    \KeywordTok{let}\NormalTok{ a}\OperatorTok{:} \DataTypeTok{Option}\OperatorTok{\textless{}}\DataTypeTok{i32}\OperatorTok{\textgreater{}} \OperatorTok{=} \ConstantTok{Some}\NormalTok{(}\DecValTok{7}\NormalTok{)}\OperatorTok{;}
    \KeywordTok{let}\NormalTok{ b}\OperatorTok{:} \DataTypeTok{Option}\OperatorTok{\textless{}}\DataTypeTok{i32}\OperatorTok{\textgreater{}} \OperatorTok{=} \ConstantTok{None}\OperatorTok{;}
    \PreprocessorTok{println!}\NormalTok{(}\StringTok{"\{\} \{\}"}\OperatorTok{,}\NormalTok{ a}\OperatorTok{.}\NormalTok{unwrap\_or(}\DecValTok{0}\NormalTok{)}\OperatorTok{,}\NormalTok{ b}\OperatorTok{.}\NormalTok{unwrap\_or(}\DecValTok{0}\NormalTok{))}\OperatorTok{;}
\OperatorTok{\}}
\end{Highlighting}
\end{Shaded}

\begin{center}\rule{0.5\linewidth}{0.5pt}\end{center}

\textbf{What is the output of this program?}

\begin{Shaded}
\begin{Highlighting}[]
\KeywordTok{fn}\NormalTok{ main() }\OperatorTok{\{}
    \KeywordTok{let}\NormalTok{ total}\OperatorTok{:} \DataTypeTok{i32} \OperatorTok{=}\NormalTok{ (}\DecValTok{1}\OperatorTok{..=}\DecValTok{5}\NormalTok{)}
        \OperatorTok{.}\NormalTok{filter(}\OperatorTok{|}\NormalTok{x}\OperatorTok{|}\NormalTok{ x }\OperatorTok{\%} \DecValTok{2} \OperatorTok{==} \DecValTok{1}\NormalTok{)}
        \OperatorTok{.}\NormalTok{map(}\OperatorTok{|}\NormalTok{x}\OperatorTok{|}\NormalTok{ x }\OperatorTok{*}\NormalTok{ x)}
        \OperatorTok{.}\NormalTok{sum()}\OperatorTok{;}
    \PreprocessorTok{println!}\NormalTok{(}\StringTok{"\{\}"}\OperatorTok{,}\NormalTok{ total)}\OperatorTok{;}
\OperatorTok{\}}
\end{Highlighting}
\end{Shaded}

\begin{center}\rule{0.5\linewidth}{0.5pt}\end{center}

\textbf{What is the output of this program?}

\begin{Shaded}
\begin{Highlighting}[]
\KeywordTok{fn}\NormalTok{ main() }\OperatorTok{\{}
    \KeywordTok{let}\NormalTok{ names }\OperatorTok{=} \PreprocessorTok{vec!}\NormalTok{[}\StringTok{"Alice"}\OperatorTok{,} \StringTok{"Bob"}\OperatorTok{,} \StringTok{"Carol"}\NormalTok{]}\OperatorTok{;}
    \KeywordTok{let}\NormalTok{ scores }\OperatorTok{=} \PreprocessorTok{vec!}\NormalTok{[}\DecValTok{90}\OperatorTok{,} \DecValTok{85}\OperatorTok{,} \DecValTok{95}\NormalTok{]}\OperatorTok{;}
    \ControlFlowTok{for}\NormalTok{ (i}\OperatorTok{,}\NormalTok{ (n}\OperatorTok{,}\NormalTok{ s)) }\KeywordTok{in}\NormalTok{ names}\OperatorTok{.}\NormalTok{iter()}\OperatorTok{.}\NormalTok{zip(scores}\OperatorTok{.}\NormalTok{iter())}\OperatorTok{.}\NormalTok{enumerate() }\OperatorTok{\{}
        \PreprocessorTok{println!}\NormalTok{(}\StringTok{"\{\}: \{\} = \{\}"}\OperatorTok{,}\NormalTok{ i}\OperatorTok{,}\NormalTok{ n}\OperatorTok{,}\NormalTok{ s)}\OperatorTok{;}
    \OperatorTok{\}}
\OperatorTok{\}}
\end{Highlighting}
\end{Shaded}

\begin{center}\rule{0.5\linewidth}{0.5pt}\end{center}

\textbf{What is the output of this program?}

\begin{Shaded}
\begin{Highlighting}[]
\KeywordTok{fn}\NormalTok{ main() }\OperatorTok{\{}
    \KeywordTok{let}\NormalTok{ greeting }\OperatorTok{=} \DataTypeTok{String}\PreprocessorTok{::}\NormalTok{from(}\StringTok{"Hello"}\NormalTok{)}\OperatorTok{;}
    \KeywordTok{let}\NormalTok{ say }\OperatorTok{=} \KeywordTok{move} \OperatorTok{||} \PreprocessorTok{println!}\NormalTok{(}\StringTok{"\{\}, world!"}\OperatorTok{,}\NormalTok{ greeting)}\OperatorTok{;}
\NormalTok{    say()}\OperatorTok{;}
\NormalTok{    say()}\OperatorTok{;}
\OperatorTok{\}}
\end{Highlighting}
\end{Shaded}

\begin{center}\rule{0.5\linewidth}{0.5pt}\end{center}

\textbf{What is the output of this program?}

\begin{Shaded}
\begin{Highlighting}[]
\AttributeTok{\#[}\NormalTok{derive}\AttributeTok{(}\BuiltInTok{Debug}\AttributeTok{)]}
\KeywordTok{struct}\NormalTok{ Point }\OperatorTok{\{}
\NormalTok{    x}\OperatorTok{:} \DataTypeTok{i32}\OperatorTok{,}
\NormalTok{    y}\OperatorTok{:} \DataTypeTok{i32}\OperatorTok{,}
\OperatorTok{\}}

\KeywordTok{fn}\NormalTok{ main() }\OperatorTok{\{}
    \KeywordTok{let}\NormalTok{ p }\OperatorTok{=}\NormalTok{ Point }\OperatorTok{\{}\NormalTok{ x}\OperatorTok{:} \DecValTok{1}\OperatorTok{,}\NormalTok{ y}\OperatorTok{:} \DecValTok{2} \OperatorTok{\};}
    \PreprocessorTok{println!}\NormalTok{(}\StringTok{"\{:?\}"}\OperatorTok{,}\NormalTok{ p)}\OperatorTok{;}
\OperatorTok{\}}
\end{Highlighting}
\end{Shaded}

\begin{center}\rule{0.5\linewidth}{0.5pt}\end{center}

\textbf{What is the output of this program?}

\begin{Shaded}
\begin{Highlighting}[]
\KeywordTok{use} \PreprocessorTok{std::collections::}\NormalTok{HashMap}\OperatorTok{;}

\KeywordTok{fn}\NormalTok{ main() }\OperatorTok{\{}
    \KeywordTok{let} \KeywordTok{mut}\NormalTok{ map}\OperatorTok{:}\NormalTok{ HashMap}\OperatorTok{\textless{}\&}\DataTypeTok{str}\OperatorTok{,} \DataTypeTok{i32}\OperatorTok{\textgreater{}} \OperatorTok{=} \PreprocessorTok{HashMap::}\NormalTok{new()}\OperatorTok{;}
\NormalTok{    map}\OperatorTok{.}\NormalTok{insert(}\StringTok{"a"}\OperatorTok{,} \DecValTok{1}\NormalTok{)}\OperatorTok{;}
    \KeywordTok{let}\NormalTok{ prev }\OperatorTok{=}\NormalTok{ map}\OperatorTok{.}\NormalTok{insert(}\StringTok{"a"}\OperatorTok{,} \DecValTok{2}\NormalTok{)}\OperatorTok{;}
    \PreprocessorTok{println!}\NormalTok{(}\StringTok{"\{:?\} \{:?\}"}\OperatorTok{,}\NormalTok{ prev}\OperatorTok{,}\NormalTok{ map}\OperatorTok{.}\NormalTok{get(}\StringTok{"a"}\NormalTok{))}\OperatorTok{;}
\OperatorTok{\}}
\end{Highlighting}
\end{Shaded}

\begin{center}\rule{0.5\linewidth}{0.5pt}\end{center}

\textbf{What is the output of this program?}

\begin{Shaded}
\begin{Highlighting}[]
\KeywordTok{use} \PreprocessorTok{std::collections::}\NormalTok{HashMap}\OperatorTok{;}

\KeywordTok{fn}\NormalTok{ main() }\OperatorTok{\{}
    \KeywordTok{let} \KeywordTok{mut}\NormalTok{ counts}\OperatorTok{:}\NormalTok{ HashMap}\OperatorTok{\textless{}}\DataTypeTok{char}\OperatorTok{,} \DataTypeTok{i32}\OperatorTok{\textgreater{}} \OperatorTok{=} \PreprocessorTok{HashMap::}\NormalTok{new()}\OperatorTok{;}
    \ControlFlowTok{for}\NormalTok{ c }\KeywordTok{in} \StringTok{"banana"}\OperatorTok{.}\NormalTok{chars() }\OperatorTok{\{}
        \OperatorTok{*}\NormalTok{counts}\OperatorTok{.}\NormalTok{entry(c)}\OperatorTok{.}\NormalTok{or\_insert(}\DecValTok{0}\NormalTok{) }\OperatorTok{+=} \DecValTok{1}\OperatorTok{;}
    \OperatorTok{\}}
    \KeywordTok{let} \KeywordTok{mut}\NormalTok{ items}\OperatorTok{:} \DataTypeTok{Vec}\OperatorTok{\textless{}}\NormalTok{(}\DataTypeTok{char}\OperatorTok{,} \DataTypeTok{i32}\NormalTok{)}\OperatorTok{\textgreater{}} \OperatorTok{=}\NormalTok{ counts}\OperatorTok{.}\NormalTok{into\_iter()}\OperatorTok{.}\NormalTok{collect()}\OperatorTok{;}
\NormalTok{    items}\OperatorTok{.}\NormalTok{sort()}\OperatorTok{;}
    \PreprocessorTok{println!}\NormalTok{(}\StringTok{"\{:?\}"}\OperatorTok{,}\NormalTok{ items)}\OperatorTok{;}
\OperatorTok{\}}
\end{Highlighting}
\end{Shaded}

\begin{center}\rule{0.5\linewidth}{0.5pt}\end{center}

\textbf{What is the output of this program?}

\begin{Shaded}
\begin{Highlighting}[]
\KeywordTok{use} \PreprocessorTok{std::collections::}\NormalTok{HashSet}\OperatorTok{;}

\KeywordTok{fn}\NormalTok{ main() }\OperatorTok{\{}
    \KeywordTok{let} \KeywordTok{mut}\NormalTok{ seen}\OperatorTok{:}\NormalTok{ HashSet}\OperatorTok{\textless{}}\DataTypeTok{i32}\OperatorTok{\textgreater{}} \OperatorTok{=} \PreprocessorTok{HashSet::}\NormalTok{new()}\OperatorTok{;}
    \PreprocessorTok{println!}\NormalTok{(}\StringTok{"\{\}"}\OperatorTok{,}\NormalTok{ seen}\OperatorTok{.}\NormalTok{insert(}\DecValTok{1}\NormalTok{))}\OperatorTok{;}
    \PreprocessorTok{println!}\NormalTok{(}\StringTok{"\{\}"}\OperatorTok{,}\NormalTok{ seen}\OperatorTok{.}\NormalTok{insert(}\DecValTok{2}\NormalTok{))}\OperatorTok{;}
    \PreprocessorTok{println!}\NormalTok{(}\StringTok{"\{\}"}\OperatorTok{,}\NormalTok{ seen}\OperatorTok{.}\NormalTok{insert(}\DecValTok{1}\NormalTok{))}\OperatorTok{;}
\OperatorTok{\}}
\end{Highlighting}
\end{Shaded}

\begin{center}\rule{0.5\linewidth}{0.5pt}\end{center}

\textbf{What is the output of this program?}

\begin{Shaded}
\begin{Highlighting}[]
\KeywordTok{use} \PreprocessorTok{std::collections::}\NormalTok{BTreeMap}\OperatorTok{;}

\KeywordTok{fn}\NormalTok{ main() }\OperatorTok{\{}
    \KeywordTok{let} \KeywordTok{mut}\NormalTok{ scores}\OperatorTok{:}\NormalTok{ BTreeMap}\OperatorTok{\textless{}\&}\DataTypeTok{str}\OperatorTok{,} \DataTypeTok{i32}\OperatorTok{\textgreater{}} \OperatorTok{=} \PreprocessorTok{BTreeMap::}\NormalTok{new()}\OperatorTok{;}
\NormalTok{    scores}\OperatorTok{.}\NormalTok{insert(}\StringTok{"carol"}\OperatorTok{,} \DecValTok{95}\NormalTok{)}\OperatorTok{;}
\NormalTok{    scores}\OperatorTok{.}\NormalTok{insert(}\StringTok{"alice"}\OperatorTok{,} \DecValTok{90}\NormalTok{)}\OperatorTok{;}
\NormalTok{    scores}\OperatorTok{.}\NormalTok{insert(}\StringTok{"bob"}\OperatorTok{,} \DecValTok{85}\NormalTok{)}\OperatorTok{;}
    \ControlFlowTok{for}\NormalTok{ (name}\OperatorTok{,}\NormalTok{ score) }\KeywordTok{in} \OperatorTok{\&}\NormalTok{scores }\OperatorTok{\{}
        \PreprocessorTok{println!}\NormalTok{(}\StringTok{"\{\} \{\}"}\OperatorTok{,}\NormalTok{ name}\OperatorTok{,}\NormalTok{ score)}\OperatorTok{;}
    \OperatorTok{\}}
\OperatorTok{\}}
\end{Highlighting}
\end{Shaded}

\begin{center}\rule{0.5\linewidth}{0.5pt}\end{center}

\textbf{What is the output of this program?}

\begin{Shaded}
\begin{Highlighting}[]
\KeywordTok{use} \PreprocessorTok{std::collections::}\NormalTok{BinaryHeap}\OperatorTok{;}

\KeywordTok{fn}\NormalTok{ main() }\OperatorTok{\{}
    \KeywordTok{let} \KeywordTok{mut}\NormalTok{ heap}\OperatorTok{:}\NormalTok{ BinaryHeap}\OperatorTok{\textless{}}\DataTypeTok{i32}\OperatorTok{\textgreater{}} \OperatorTok{=} \PreprocessorTok{BinaryHeap::}\NormalTok{new()}\OperatorTok{;}
    \ControlFlowTok{for}\NormalTok{ x }\KeywordTok{in}\NormalTok{ [}\DecValTok{3}\OperatorTok{,} \DecValTok{1}\OperatorTok{,} \DecValTok{4}\OperatorTok{,} \DecValTok{1}\OperatorTok{,} \DecValTok{5}\OperatorTok{,} \DecValTok{9}\OperatorTok{,} \DecValTok{2}\NormalTok{] }\OperatorTok{\{}
\NormalTok{        heap}\OperatorTok{.}\NormalTok{push(x)}\OperatorTok{;}
    \OperatorTok{\}}
    \ControlFlowTok{while} \KeywordTok{let} \ConstantTok{Some}\NormalTok{(top) }\OperatorTok{=}\NormalTok{ heap}\OperatorTok{.}\NormalTok{pop() }\OperatorTok{\{}
        \PreprocessorTok{print!}\NormalTok{(}\StringTok{"\{\} "}\OperatorTok{,}\NormalTok{ top)}\OperatorTok{;}
    \OperatorTok{\}}
    \PreprocessorTok{println!}\NormalTok{()}\OperatorTok{;}
\OperatorTok{\}}
\end{Highlighting}
\end{Shaded}

\begin{center}\rule{0.5\linewidth}{0.5pt}\end{center}

\textbf{What is the output of this program?}

\begin{Shaded}
\begin{Highlighting}[]
\KeywordTok{use} \PreprocessorTok{std::collections::}\NormalTok{VecDeque}\OperatorTok{;}

\KeywordTok{fn}\NormalTok{ main() }\OperatorTok{\{}
    \KeywordTok{let} \KeywordTok{mut}\NormalTok{ q}\OperatorTok{:}\NormalTok{ VecDeque}\OperatorTok{\textless{}}\DataTypeTok{i32}\OperatorTok{\textgreater{}} \OperatorTok{=} \PreprocessorTok{VecDeque::}\NormalTok{new()}\OperatorTok{;}
\NormalTok{    q}\OperatorTok{.}\NormalTok{push\_back(}\DecValTok{1}\NormalTok{)}\OperatorTok{;}
\NormalTok{    q}\OperatorTok{.}\NormalTok{push\_back(}\DecValTok{2}\NormalTok{)}\OperatorTok{;}
\NormalTok{    q}\OperatorTok{.}\NormalTok{push\_front(}\DecValTok{0}\NormalTok{)}\OperatorTok{;}
\NormalTok{    q}\OperatorTok{.}\NormalTok{push\_back(}\DecValTok{3}\NormalTok{)}\OperatorTok{;}
    \ControlFlowTok{while} \KeywordTok{let} \ConstantTok{Some}\NormalTok{(x) }\OperatorTok{=}\NormalTok{ q}\OperatorTok{.}\NormalTok{pop\_front() }\OperatorTok{\{}
        \PreprocessorTok{print!}\NormalTok{(}\StringTok{"\{\} "}\OperatorTok{,}\NormalTok{ x)}\OperatorTok{;}
    \OperatorTok{\}}
    \PreprocessorTok{println!}\NormalTok{()}\OperatorTok{;}
\OperatorTok{\}}
\end{Highlighting}
\end{Shaded}

\begin{center}\rule{0.5\linewidth}{0.5pt}\end{center}

\textbf{What is the output of this program?}

\begin{Shaded}
\begin{Highlighting}[]
\KeywordTok{fn}\NormalTok{ classify(x}\OperatorTok{:} \DataTypeTok{i32}\NormalTok{) }\OperatorTok{{-}\textgreater{}} \OperatorTok{\&}\OtherTok{\textquotesingle{}static} \DataTypeTok{str} \OperatorTok{\{}
    \ControlFlowTok{match}\NormalTok{ x }\OperatorTok{\{}
\NormalTok{        n }\ControlFlowTok{if}\NormalTok{ n }\OperatorTok{\textless{}} \DecValTok{0} \OperatorTok{=\textgreater{}} \StringTok{"negative"}\OperatorTok{,}
        \DecValTok{0} \OperatorTok{=\textgreater{}} \StringTok{"zero"}\OperatorTok{,}
        \DecValTok{1}\OperatorTok{..=}\DecValTok{9} \OperatorTok{=\textgreater{}} \StringTok{"single digit"}\OperatorTok{,}
\NormalTok{        \_ }\OperatorTok{=\textgreater{}} \StringTok{"big"}\OperatorTok{,}
    \OperatorTok{\}}
\OperatorTok{\}}

\KeywordTok{fn}\NormalTok{ main() }\OperatorTok{\{}
    \ControlFlowTok{for}\NormalTok{ x }\KeywordTok{in}\NormalTok{ [}\OperatorTok{{-}}\DecValTok{1}\OperatorTok{,} \DecValTok{0}\OperatorTok{,} \DecValTok{7}\OperatorTok{,} \DecValTok{42}\NormalTok{] }\OperatorTok{\{}
        \PreprocessorTok{println!}\NormalTok{(}\StringTok{"\{\}: \{\}"}\OperatorTok{,}\NormalTok{ x}\OperatorTok{,}\NormalTok{ classify(x))}\OperatorTok{;}
    \OperatorTok{\}}
\OperatorTok{\}}
\end{Highlighting}
\end{Shaded}

\begin{center}\rule{0.5\linewidth}{0.5pt}\end{center}

\textbf{What is the output of this program?}

\begin{Shaded}
\begin{Highlighting}[]
\KeywordTok{fn}\NormalTok{ longest}\OperatorTok{\textless{}}\OtherTok{\textquotesingle{}a}\OperatorTok{\textgreater{}}\NormalTok{(x}\OperatorTok{:} \OperatorTok{\&}\OtherTok{\textquotesingle{}a} \DataTypeTok{str}\OperatorTok{,}\NormalTok{ y}\OperatorTok{:} \OperatorTok{\&}\OtherTok{\textquotesingle{}a} \DataTypeTok{str}\NormalTok{) }\OperatorTok{{-}\textgreater{}} \OperatorTok{\&}\OtherTok{\textquotesingle{}a} \DataTypeTok{str} \OperatorTok{\{}
    \ControlFlowTok{if}\NormalTok{ x}\OperatorTok{.}\NormalTok{len() }\OperatorTok{\textgreater{}=}\NormalTok{ y}\OperatorTok{.}\NormalTok{len() }\OperatorTok{\{}\NormalTok{ x }\OperatorTok{\}} \ControlFlowTok{else} \OperatorTok{\{}\NormalTok{ y }\OperatorTok{\}}
\OperatorTok{\}}

\KeywordTok{fn}\NormalTok{ main() }\OperatorTok{\{}
    \KeywordTok{let}\NormalTok{ s1 }\OperatorTok{=} \DataTypeTok{String}\PreprocessorTok{::}\NormalTok{from(}\StringTok{"hello"}\NormalTok{)}\OperatorTok{;}
    \KeywordTok{let}\NormalTok{ s2 }\OperatorTok{=} \DataTypeTok{String}\PreprocessorTok{::}\NormalTok{from(}\StringTok{"ds210"}\NormalTok{)}\OperatorTok{;}
    \KeywordTok{let}\NormalTok{ result }\OperatorTok{=}\NormalTok{ longest(}\OperatorTok{\&}\NormalTok{s1}\OperatorTok{,} \OperatorTok{\&}\NormalTok{s2)}\OperatorTok{;}
    \PreprocessorTok{println!}\NormalTok{(}\StringTok{"\{\}"}\OperatorTok{,}\NormalTok{ result)}\OperatorTok{;}
\OperatorTok{\}}
\end{Highlighting}
\end{Shaded}

\begin{center}\rule{0.5\linewidth}{0.5pt}\end{center}

\textbf{What is the output of this program?}

\begin{Shaded}
\begin{Highlighting}[]
\KeywordTok{struct}\NormalTok{ Rectangle }\OperatorTok{\{}
\NormalTok{    width}\OperatorTok{:} \DataTypeTok{u32}\OperatorTok{,}
\NormalTok{    height}\OperatorTok{:} \DataTypeTok{u32}\OperatorTok{,}
\OperatorTok{\}}

\KeywordTok{impl}\NormalTok{ Rectangle }\OperatorTok{\{}
    \KeywordTok{fn}\NormalTok{ area(}\OperatorTok{\&}\KeywordTok{self}\NormalTok{) }\OperatorTok{{-}\textgreater{}} \DataTypeTok{u32} \OperatorTok{\{}
        \KeywordTok{self}\OperatorTok{.}\NormalTok{width }\OperatorTok{*} \KeywordTok{self}\OperatorTok{.}\NormalTok{height}
    \OperatorTok{\}}
\OperatorTok{\}}

\KeywordTok{fn}\NormalTok{ main() }\OperatorTok{\{}
    \KeywordTok{let}\NormalTok{ r }\OperatorTok{=}\NormalTok{ Rectangle }\OperatorTok{\{}\NormalTok{ width}\OperatorTok{:} \DecValTok{3}\OperatorTok{,}\NormalTok{ height}\OperatorTok{:} \DecValTok{4} \OperatorTok{\};}
    \PreprocessorTok{println!}\NormalTok{(}\StringTok{"\{\}"}\OperatorTok{,}\NormalTok{ r}\OperatorTok{.}\NormalTok{area())}\OperatorTok{;}
\OperatorTok{\}}
\end{Highlighting}
\end{Shaded}

\begin{center}\rule{0.5\linewidth}{0.5pt}\end{center}

\subsection{Short Coding Challenges (coding + short
tag)}\label{short-coding-challenges-coding-short-tag}

\textbf{Write a function
\texttt{char\_count(s:\ \&str)\ -\textgreater{}\ HashMap\textless{}char,\ usize\textgreater{}}
that returns a map from each character in \texttt{s} to the number of
times it appears.}

Call it in \texttt{main} with \texttt{"rust"} and print the resulting
map with \texttt{\{:?\}}.

\begin{Shaded}
\begin{Highlighting}[]


\end{Highlighting}
\end{Shaded}

\begin{center}\rule{0.5\linewidth}{0.5pt}\end{center}

\textbf{In \texttt{main}, start with
\texttt{let\ nums\ =\ vec!{[}1,\ 2,\ 2,\ 3,\ 1,\ 4,\ 3,\ 5{]};} and use
a \texttt{HashSet\textless{}i32\textgreater{}} to count how many
distinct values appear. Print the count.}

\begin{Shaded}
\begin{Highlighting}[]


\end{Highlighting}
\end{Shaded}

\begin{center}\rule{0.5\linewidth}{0.5pt}\end{center}

\textbf{In \texttt{main}, build a
\texttt{BTreeMap\textless{}\&str,\ i32\textgreater{}} from the pairs
\texttt{("banana",\ 3)}, \texttt{("apple",\ 5)},
\texttt{("cherry",\ 2)}, and then iterate and print each
\texttt{name\ score} pair on its own line.}

The output should be in alphabetical order by name because of the map
type.

\begin{Shaded}
\begin{Highlighting}[]


\end{Highlighting}
\end{Shaded}

\begin{center}\rule{0.5\linewidth}{0.5pt}\end{center}

\textbf{In \texttt{main}, use a
\texttt{BinaryHeap\textless{}i32\textgreater{}} to find and print the
three largest values in \texttt{vec!{[}4,\ 10,\ 2,\ 8,\ 15,\ 7,\ 1{]}},
in descending order, separated by spaces.}

\begin{Shaded}
\begin{Highlighting}[]


\end{Highlighting}
\end{Shaded}

\begin{center}\rule{0.5\linewidth}{0.5pt}\end{center}

\textbf{In \texttt{main}, use
\texttt{VecDeque\textless{}i32\textgreater{}} as a FIFO queue. Enqueue
the numbers 1 through 5 (in order), then dequeue them all and print each
dequeued value on its own line.}

\begin{Shaded}
\begin{Highlighting}[]


\end{Highlighting}
\end{Shaded}

\begin{center}\rule{0.5\linewidth}{0.5pt}\end{center}

\textbf{In \texttt{main}, use iterator methods to compute and print the
sum of squares of the even numbers in \texttt{1..=10}.}

\begin{Shaded}
\begin{Highlighting}[]


\end{Highlighting}
\end{Shaded}

\begin{center}\rule{0.5\linewidth}{0.5pt}\end{center}

\textbf{In \texttt{main}, given
\texttt{vec!{[}3,\ 1,\ 4,\ 1,\ 5,\ 9,\ 2,\ 6,\ 5{]}}, use
\texttt{.partition(...)} to split the values into two
\texttt{Vec\textless{}i32\textgreater{}}s: those less than 5 and those
greater than or equal to 5. Then print the lengths of the two resulting
vectors.}

\begin{Shaded}
\begin{Highlighting}[]


\end{Highlighting}
\end{Shaded}

\begin{center}\rule{0.5\linewidth}{0.5pt}\end{center}

\textbf{Given
\texttt{let\ fruits\ =\ vec!{[}"apple",\ "banana",\ "apple",\ "cherry",\ "banana",\ "apple"{]};},
in \texttt{main} build a
\texttt{HashMap\textless{}\&str,\ i32\textgreater{}} using the
\texttt{entry} API that inserts each key and increments the count for
each key and print how many times \texttt{"apple"} appears.}

\begin{Shaded}
\begin{Highlighting}[]


\end{Highlighting}
\end{Shaded}

\begin{center}\rule{0.5\linewidth}{0.5pt}\end{center}

\textbf{Write a function
\texttt{first\_long\textless{}\textquotesingle{}a\textgreater{}(words:\ \&\textquotesingle{}a\ {[}\&\textquotesingle{}a\ str{]},\ min\_len:\ usize)\ -\textgreater{}\ Option\textless{}\&\textquotesingle{}a\ str\textgreater{}}
that returns the first word in \texttt{words} whose length is at least
\texttt{min\_len}, or \texttt{None} if there isn't one.}

Call it from \texttt{main} with
\texttt{\&{[}"hi",\ "ok",\ "rust",\ "hello"{]}} and
\texttt{min\_len\ =\ 4}, and print the result using \texttt{\{:?\}}.

\begin{Shaded}
\begin{Highlighting}[]


\end{Highlighting}
\end{Shaded}

\begin{center}\rule{0.5\linewidth}{0.5pt}\end{center}

\textbf{Write a generic function
\texttt{largest\textless{}T:\ PartialOrd\ +\ Copy\textgreater{}(list:\ \&{[}T{]})\ -\textgreater{}\ T}
that returns the largest element of the slice. Assume the slice is
non-empty.}

Call it from \texttt{main} once with
\texttt{\&{[}10,\ 3,\ 7,\ 25,\ 4{]}} and once with
\texttt{\&{[}1.5\_f64,\ 2.5,\ 0.5{]}} and print both results.

\begin{Shaded}
\begin{Highlighting}[]


\end{Highlighting}
\end{Shaded}

\begin{center}\rule{0.5\linewidth}{0.5pt}\end{center}

\textbf{In \texttt{main}, start with
\texttt{let\ mut\ scores\ =\ vec!{[}("alice",\ 85),\ ("bob",\ 92),\ ("carol",\ 78){]};}
and use \texttt{sort\_by} with a closure to sort the vector by score in
descending order. Print the resulting vector with \texttt{\{:?\}}.}

\begin{Shaded}
\begin{Highlighting}[]


\end{Highlighting}
\end{Shaded}

\begin{center}\rule{0.5\linewidth}{0.5pt}\end{center}

\subsection{Long Coding Challenges (coding + long
tag)}\label{long-coding-challenges-coding-long-tag}

\textbf{Write a program that counts word frequencies in a sentence and
prints the three most frequent words, along with their counts, in
descending order of frequency.}

Use the sentence
\texttt{"the\ quick\ brown\ fox\ jumps\ over\ the\ lazy\ dog\ the\ fox\ runs\ fast\ over\ the\ hill"}.
Split on whitespace to get words (do not worry about punctuation or
case). Break ties by alphabetical order of the word.

Expected output:

\begin{Shaded}
\begin{Highlighting}[]
\NormalTok{the 4}
\NormalTok{fox 2}
\NormalTok{over 2}
\end{Highlighting}
\end{Shaded}

\begin{Shaded}
\begin{Highlighting}[]


\end{Highlighting}
\end{Shaded}

\begin{center}\rule{0.5\linewidth}{0.5pt}\end{center}

\textbf{Write a program that prints, in sorted order, the values that
appear in both input vectors:}

\begin{Shaded}
\begin{Highlighting}[]
\KeywordTok{let}\NormalTok{ a }\OperatorTok{=} \PreprocessorTok{vec!}\NormalTok{[}\DecValTok{1}\OperatorTok{,} \DecValTok{2}\OperatorTok{,} \DecValTok{3}\OperatorTok{,} \DecValTok{4}\OperatorTok{,} \DecValTok{5}\OperatorTok{,} \DecValTok{6}\NormalTok{]}\OperatorTok{;}
\KeywordTok{let}\NormalTok{ b }\OperatorTok{=} \PreprocessorTok{vec!}\NormalTok{[}\DecValTok{4}\OperatorTok{,} \DecValTok{5}\OperatorTok{,} \DecValTok{6}\OperatorTok{,} \DecValTok{7}\OperatorTok{,} \DecValTok{8}\OperatorTok{,} \DecValTok{9}\NormalTok{]}\OperatorTok{;}
\end{Highlighting}
\end{Shaded}

\textbf{Use \texttt{HashSet} for the membership checks, then sort the
resulting common values before printing them (one per line).}

Expected output:

\begin{Shaded}
\begin{Highlighting}[]
\NormalTok{4}
\NormalTok{5}
\NormalTok{6}
\end{Highlighting}
\end{Shaded}

\begin{Shaded}
\begin{Highlighting}[]


\end{Highlighting}
\end{Shaded}

\begin{center}\rule{0.5\linewidth}{0.5pt}\end{center}

\textbf{Build a small ``grade book''. Given
\texttt{let\ entries\ =\ vec!{[}("carol",\ 72),\ ("alice",\ 95),\ ("bob",\ 58),\ ("alice",\ 88),\ ("bob",\ 80){]};},
insert each pair into a
\texttt{BTreeMap\textless{}\&str,\ Vec\textless{}i32\textgreater{}\textgreater{}}
keyed by name, with the list of that student's scores as the value. Then
iterate the map (which will give names in alphabetical order) and print
each student's name followed by their average score, formatted to one
decimal place.}

Expected output:

\begin{Shaded}
\begin{Highlighting}[]
\NormalTok{alice 91.5}
\NormalTok{bob 69.0}
\NormalTok{carol 72.0}
\end{Highlighting}
\end{Shaded}

\begin{Shaded}
\begin{Highlighting}[]


\end{Highlighting}
\end{Shaded}

\begin{center}\rule{0.5\linewidth}{0.5pt}\end{center}

\textbf{Write a function
\texttt{kth\_largest(nums:\ \&{[}i32{]},\ k:\ usize)\ -\textgreater{}\ Option\textless{}i32\textgreater{}}
that returns the \texttt{k}-th largest value in the slice (1-indexed).
If \texttt{k} is zero or larger than the slice length, return
\texttt{None}. Use a \texttt{BinaryHeap} to find the answer.}

Call the function from \texttt{main} with
\texttt{\&{[}7,\ 2,\ 9,\ 4,\ 11,\ 3{]}} and \texttt{k\ =\ 3}, and print
the result with \texttt{\{:?\}}.

Expected output:

\begin{Shaded}
\begin{Highlighting}[]
\NormalTok{Some(7)}
\end{Highlighting}
\end{Shaded}

\begin{Shaded}
\begin{Highlighting}[]


\end{Highlighting}
\end{Shaded}

\begin{center}\rule{0.5\linewidth}{0.5pt}\end{center}

\textbf{Write a function
\texttt{is\_balanced(s:\ \&str)\ -\textgreater{}\ bool} that returns
\texttt{true} if every opening bracket in \texttt{s} is matched by a
closing bracket of the same type in the correct order. Support
\texttt{()}, \texttt{{[}{]}}, and \texttt{\{\}}. Ignore any other
characters.}

Use a \texttt{Vec\textless{}char\textgreater{}} as a stack. Call the
function from \texttt{main} on each of these inputs and print the result
on its own line:

\begin{Shaded}
\begin{Highlighting}[]
\KeywordTok{let}\NormalTok{ cases }\OperatorTok{=}\NormalTok{ [}\StringTok{"()[]\{\}"}\OperatorTok{,} \StringTok{"([\{\}])"}\OperatorTok{,} \StringTok{"(]"}\OperatorTok{,} \StringTok{"((()"}\NormalTok{]}\OperatorTok{;}
\end{Highlighting}
\end{Shaded}

Expected output:

\begin{Shaded}
\begin{Highlighting}[]
\NormalTok{true}
\NormalTok{true}
\NormalTok{false}
\NormalTok{false}
\end{Highlighting}
\end{Shaded}

\begin{Shaded}
\begin{Highlighting}[]


\end{Highlighting}
\end{Shaded}

\begin{center}\rule{0.5\linewidth}{0.5pt}\end{center}

\textbf{Write a function
\texttt{moving\_avg(values:\ \&{[}f64{]},\ window:\ usize)\ -\textgreater{}\ Vec\textless{}f64\textgreater{}}
that computes the rolling average over a sliding window of size
\texttt{window}. Use a \texttt{VecDeque\textless{}f64\textgreater{}} to
track the window. Only emit an average once the window is fully
populated, so the output length is
\texttt{values.len()\ -\ window\ +\ 1} (or empty if
\texttt{values.len()\ \textless{}\ window}).}

Call the function from \texttt{main} with
\texttt{values\ =\ \&{[}1.0,\ 2.0,\ 3.0,\ 4.0,\ 5.0{]}} and
\texttt{window\ =\ 3}, and print each result on its own line, formatted
to one decimal place.

Expected output:

\begin{Shaded}
\begin{Highlighting}[]
\NormalTok{2.0}
\NormalTok{3.0}
\NormalTok{4.0}
\end{Highlighting}
\end{Shaded}

\begin{Shaded}
\begin{Highlighting}[]


\end{Highlighting}
\end{Shaded}

\begin{center}\rule{0.5\linewidth}{0.5pt}\end{center}

\textbf{Define a trait \texttt{Summary} with a single method
\texttt{summary(\&self)\ -\textgreater{}\ String}. Implement
\texttt{Summary} for two structs:
\texttt{Tweet\ \{\ user:\ String,\ text:\ String\ \}} and
\texttt{Article\ \{\ title:\ String,\ body:\ String\ \}}. The tweet
summary should be \texttt{"@user:\ text"}, and the article summary
should be \texttt{"title\ —\ first\ 20\ chars\ of\ body"} (use
\texttt{.chars().take(20).collect::\textless{}String\textgreater{}()} to
get the first 20 characters).}

In \texttt{main}, create one of each, put both into a
\texttt{Vec\textless{}Box\textless{}dyn\ Summary\textgreater{}\textgreater{}},
iterate over the vector, and print each summary.

\begin{Shaded}
\begin{Highlighting}[]


\end{Highlighting}
\end{Shaded}

\begin{center}\rule{0.5\linewidth}{0.5pt}\end{center}

\textbf{Given the slice of \texttt{(city,\ temp\_f)} pairs:}

\begin{Shaded}
\begin{Highlighting}[]
\KeywordTok{let}\NormalTok{ readings }\OperatorTok{=} \PreprocessorTok{vec!}\NormalTok{[}
\NormalTok{    (}\StringTok{"Boston"}\OperatorTok{,} \DecValTok{48.0\_f64}\NormalTok{)}\OperatorTok{,}\NormalTok{ (}\StringTok{"Phoenix"}\OperatorTok{,} \DecValTok{102.5}\NormalTok{)}\OperatorTok{,}
\NormalTok{    (}\StringTok{"Denver"}\OperatorTok{,} \DecValTok{77.0}\NormalTok{)}\OperatorTok{,}\NormalTok{ (}\StringTok{"Miami"}\OperatorTok{,} \DecValTok{88.0}\NormalTok{)}\OperatorTok{,}\NormalTok{ (}\StringTok{"Anchorage"}\OperatorTok{,} \DecValTok{36.0}\NormalTok{)}\OperatorTok{,}
\NormalTok{]}\OperatorTok{;}
\end{Highlighting}
\end{Shaded}

\textbf{Write a program that prints every city whose temperature is
above 75°F, with its temperature converted to Celsius (formula:
\texttt{(f\ -\ 32.0)\ *\ 5.0\ /\ 9.0}), one per line. Sort the output
alphabetically by city. Use iterator methods (\texttt{filter},
\texttt{map}, \texttt{collect}, \texttt{sort\_by\_key}, \ldots) rather
than explicit loops.}

Each printed line should look like \texttt{Phoenix\ 39.2} (temperature
formatted to one decimal place).

\begin{Shaded}
\begin{Highlighting}[]


\end{Highlighting}
\end{Shaded}

\begin{center}\rule{0.5\linewidth}{0.5pt}\end{center}

\textbf{Define a struct
\texttt{Rectangle\ \{\ width:\ u32,\ height:\ u32\ \}} with three
methods:}

\begin{enumerate}
\def\labelenumi{\arabic{enumi}.}
\tightlist
\item
  \texttt{area(\&self)\ -\textgreater{}\ u32}
\item
  \texttt{perimeter(\&self)\ -\textgreater{}\ u32}
\item
  \texttt{can\_hold(\&self,\ other:\ \&Rectangle)\ -\textgreater{}\ bool}
  --- returns \texttt{true} if \texttt{self} is strictly larger than
  \texttt{other} in both dimensions
\end{enumerate}

\textbf{In \texttt{main}, create
\texttt{let\ a\ =\ Rectangle\ \{\ width:\ 10,\ height:\ 5\ \};} and
\texttt{let\ b\ =\ Rectangle\ \{\ width:\ 4,\ height:\ 3\ \};}. Print
\texttt{a.area()}, \texttt{a.perimeter()}, and \texttt{a.can\_hold(\&b)}
each on their own line.}

\begin{Shaded}
\begin{Highlighting}[]


\end{Highlighting}
\end{Shaded}

\begin{center}\rule{0.5\linewidth}{0.5pt}\end{center}

\textbf{You are given
\texttt{let\ scores\ =\ vec!{[}("alice",\ 90),\ ("bob",\ 75),\ ("alice",\ 82),\ ("carol",\ 88),\ ("bob",\ 91),\ ("alice",\ 79){]};}.
Using the \texttt{Entry} API on a
\texttt{HashMap\textless{}\&str,\ (i32,\ i32)\textgreater{}} (where the
value is \texttt{(total\_points,\ count)}), compute each student's
average. Then print a line for each student in alphabetical order,
formatted \texttt{name\ avg} with the average rounded to one decimal
place.}

Expected output:

\begin{Shaded}
\begin{Highlighting}[]
\NormalTok{alice 83.7}
\NormalTok{bob 83.0}
\NormalTok{carol 88.0}
\end{Highlighting}
\end{Shaded}

\begin{Shaded}
\begin{Highlighting}[]


\end{Highlighting}
\end{Shaded}

\begin{center}\rule{0.5\linewidth}{0.5pt}\end{center}

\end{document}
